\chapter{Proofs for Huber Joint Boundaries and the Sign Bridge}
\label{annex:p15-huber-joint-endpoints}

\begin{annexitem}{annex:p15.huber.joint.identified}{thm:p15-huber-joint-identified}{The exact centered scale condition and eventual unique estimation}
\begin{proof}
The function \(s\mapsto E\min(X^2/s^2,k^2)\) is continuous and
nonincreasing, by bounded convergence. Its limits at zero and
infinity are \(k^2(1-m_0)\) and zero. If the upper limit exceeds
\(\beta\), the intermediate value theorem supplies a positive root.
At such a root there is positive probability of
\(0<|X|<ks\); otherwise the moment would still equal its upper
limit. Those observations make the moment strictly decrease near
the root, proving uniqueness. At this root define
\(U=X/S_0\), \(A=P(|U|<k)\) and
\(C=E[U^2{\bf1}_{\{|U|<k\}}]\).
The same observations prove \(C>0\) and \(A>0\).
If the upper limit is below \(\beta\), no root exists.
If it equals \(\beta\), equality at any positive \(s\) requires
every nonzero observation to be saturated almost surely.
Every smaller positive scale is then also a root. This proves the
complete scale condition and the alternatives at equality.
The two examples follow by saturating the nonzero atoms, or by
noting that the uniform component always gives positive mass to
the unsaturated interval.

We prove the joint and sample assertions by a convex criterion,
so no population second moment is hidden. The perspective
\(s\rho((x-t)/s)\) is convex in \((t,s)\), \(s>0\).
Indeed at a convex combination of two parameter points, its
residual is the convex combination of the two standardized
residuals with weights proportional to their positive scales;
the defining inequality for convex \(\rho\), multiplied by the
combined scale, is exactly the required inequality.
Direct differentiation, including the matching values at the
Huber joins, gives
\[
 \nabla_{t,s}\{s\rho((x-t)/s)+\beta s/2\}
 =\begin{pmatrix}
 -\psi((x-t)/s)\\
 \{\beta-\psi((x-t)/s)^2\}/2
 \end{pmatrix}.
\]
Both components are bounded uniformly in \(x,t,s\).
Criterion differences from a fixed reference point are consequently
bounded in absolute value by
\(k|t-t_0|+(k^2+\beta)|s-s_0|/2\).
Their expectation is finite even if \(E|X|=\infty\).
Dominated convergence gives the displayed population gradient.
At \((0,S_0)\) it is zero, by symmetry and the scale equation;
convexity makes this a global minimum.

For strict local curvature choose \(\delta>0\) so that
\[
 A_\delta=P(|U|\leq k-\delta)>0,\qquad
 C_\delta=E[U^2{\bf1}_{\{|U|\leq k-\delta\}}]>0.
\]
Such a choice exists because \(C>0\). Symmetry gives zero
off-diagonal moment in the matrix
\[
 E\left[{\bf1}_{\{|U|\leq k-\delta\}}
       \begin{pmatrix}1&U\\U&U^2\end{pmatrix}\right]
       =\begin{pmatrix}A_\delta&0\\0&C_\delta\end{pmatrix}.
\]
In a sufficiently small ball about \((0,S_0)\), these fixed
observations remain strictly interior. Their individual Hessian
is \(s^{-1}(1,v)^{\mathsf T}(1,v)\), \(v=(x-t)/s\).
On this restricted set its change from the reference Hessian
is uniformly small, since \(U\) is bounded and \(s\) stays away
from zero. Thus this part of the population criterion has a
uniform positive-definite lower curvature bound. All remaining
summands are convex and cannot reduce that bound.
The minimum is locally strict. A second global minimum of a
convex function would make the joining segment constant, including
points inside this strict neighborhood, a contradiction.
This proves population uniqueness.

The same argument applies almost surely to empirical criteria:
the restricted bounded Gram matrix converges by the strong law
to the positive-definite matrix above. Hence for some fixed
neighborhood and \(\lambda>0\), eventually
\[
 (\nabla L_n(\theta)-\nabla L_n(\theta_0))
        \mathbin{\cdot}(\theta-\theta_0)
       \geq\lambda\|\theta-\theta_0\|^2
\]
there, where \(\theta_0=(0,S_0)\).
The empirical gradient at \(\theta_0\) tends almost surely to zero.
Minimize the continuous finite criterion on a small closed ball
inside the neighborhood. On its boundary the derivative toward
the center is negative once
\(\|\nabla L_n(\theta_0)\|<\lambda\) times the radius.
The minimum is therefore interior and has gradient zero.
Local strict convexity and global convexity give its uniqueness
on the entire positive-scale domain.
At that root,
\[
             \|\widehat\theta_n-\theta_0\|
              \leq\lambda^{-1}\|\nabla L_n(\theta_0)\|.
\]
It follows that the roots converge almost surely.

Measurability can be made explicit using rational parameter grids.
The infimum over rational positive-scale points equals the
criterion infimum by continuity and is measurable in the data.
Select the first rational point in a fixed enumeration within
\(1/j\) of that infimum. At an attained minimum with the proved
strict local curvature, convexity gives a positive gap outside
each sufficiently small ball around the minimum; hence these
selected points converge to it. Their limit is measurable.
Take the limit when the rational approximants converge to a finite
point of positive scale, and use \((0,1)\) otherwise.
Convergence of a sequence of measurable coordinates and positivity
of the limiting scale are measurable conditions.
On the event proved above, eventually this selection is the unique
positive root. Its other values do not affect the asymptotic claims.
No equivariance property of that exceptional extension is needed.
\end{proof}
\end{annexitem}

\begin{annexitem}{annex:p15.huber.joint.directional}{thm:p15-huber-joint-directional}{Exact population derivatives at arbitrary clipping atoms}
\begin{proof}
For a direction \(h=(a,b)\), the perturbed residual is
\[
 \frac{U-ta}{1+tb}=U-t(a+bU)+O(t^2)
\]
uniformly in bounded \(U,h\). At a strictly interior observation
the derivative of the first score is \(-a-bU\), and that of
the half-squared score is \(-U(a+bU)\).
Their interior expectations are \(-Aa,-Cb\), because the two
odd interior moments vanish by symmetry.
Strictly exterior observations have zero derivatives.
At \(U=k\), the first derivative is
\(-(a+kb)_+\); at \(U=-k\), it is \((-a+kb)_+\).
Multiplication by the unperturbed scores \(k,-k\) gives the
half-squared-score derivatives at the two atoms.
Combining their equal masses \(p\) gives exactly
\(-\mathcal D_1(h),-\mathcal D_2(h)\).

We justify the uniformity and the absence of an interchange problem.
On any compact direction set, both scores change by at most
a constant times \(t\); a residual far outside the clipping
region stays saturated. After separating the atoms at \(\pm k\),
the only unaccounted changes of branch occur in
\[
             0<\bigl||U|-k\bigr|\leq Kt
\]
for a constant \(K\) depending on that direction set.
The probability of this set tends to zero. Outside it the
residual expansion and the ordinary polynomial score expansion
have a uniform \(O(t^2)\) remainder.
At the two atoms, positive-part Lipschitz continuity gives
a uniform \(O(t^2)\) remainder as well.
After division by \(t\), the integrated remainder is bounded
by \(O(t)\) plus a constant times the displayed probability.
This tends to zero uniformly, and also permits convergent
varying directions. For \(p>0\), the derivatives along
\((0,1)\) and \((0,-1)\) have scale components
\(C+2pk^2\) and \(-C\); they cannot be values of one linear map.
\end{proof}
\end{annexitem}

\begin{annexitem}{annex:p15.huber.joint.cone.inverse}{thm:p15-huber-joint-cone-inverse}{An everywhere defined inverse with four explicit pieces}
\begin{proof}
Let \(\lambda=\min(A,C)>0\). For any \(h,g\),
\[
 (\mathcal D(h)-\mathcal D(g))\cdot(h-g)
 \geq\lambda\|h-g\|^2.
\]
The diagonal part supplies this lower bound; each added term is
\[
 p\{(v\cdot h)_+-(v\cdot g)_+\}\{v\cdot h-v\cdot g\}\geq0.
\]
Thus the map is one-to-one. For each \(z\), the continuous
strictly convex function \(\Gamma(h)-z\cdot h\) tends to infinity
with \(\|h\|\), since \(\Gamma(h)\geq\lambda\|h\|^2/2\).
It attains a unique minimum, whose gradient equation is
\(\mathcal D(h)=z\). This proves surjectivity without a
topological inverse assertion.
The displayed monotonicity inequality and Cauchy--Schwarz give
\(\|h-g\|\leq\lambda^{-1}\|\mathcal D(h)-\mathcal D(g)\|\),
proving the inverse bound.

Within a fixed pair of sign cones the positive parts select
exactly the terms in \(M_e\), so \(\mathcal D(h)=M_eh\).
Every such matrix is positive definite by its diagonal lower
bound. On a boundary \(v\cdot h=0\), adding or dropping
\(pvv^{\mathsf T}h\) changes nothing. Hence the finite procedure
has a retained piece for every \(z\), and all retained boundary
values agree by the already proved uniqueness.
\end{proof}
\end{annexitem}

\begin{annexitem}{annex:p15.huber.joint.atomic.limit}{thm:p15-huber-joint-atomic-limit}{The full joint limit, including the atoms}
\begin{proof}
The strong local curvature argument in the identification proof
already bounds the parameter error by the norm of the empirical
gradient at the true pair. The latter is an average of bounded
centered random vectors. Chebyshev's inequality therefore gives
\(\sqrt n(\widehat\theta_n-\theta_0)=O_P(1)\).
In particular \(H_n=O_P(1)\); this rate has not been assumed.

Let \(m_n(h)\) denote the empirical two-moment vector in the
directional theorem. The same remainder decomposition, now at
\(t=n^{-1/2}\), gives, uniformly on each compact set of \(h\),
\[
 \sqrt n\{m_n(h/\sqrt n)-m_n(0)\}
                         =-\mathcal D(h)+o_P(1).
\]
Here are the empirical details. The bounded strictly interior
first and second moments converge by the strong law to their
expectations. The two exact atom frequencies converge to \(p\).
The remainder is bounded by \(O(n^{-1/2})\) plus a constant
times the empirical mass of
\(0<||U|-k|\leq K/\sqrt n\).
That empirical mass tends to zero in probability by Markov's
inequality, because its expectation tends to zero. Equivalently,
fixed shrinking bands and the strong law give the same result.
These bounds hold uniformly over the chosen compact direction set.

Put
\[
 Z_n=\frac1{\sqrt n}\sum_i
       \begin{pmatrix}
       \psi(U_i)\\(\psi(U_i)^2-\beta)/2
       \end{pmatrix}.
\]
Write \(W\) for one of these centered bounded summands.
Their covariance matrix is
\(\operatorname{diag}(\beta,\eta^2)\), since the cross moment
is odd under the symmetric law.
For completeness, their bivariate Gaussian convergence follows
directly by expanding the characteristic function at any
fixed \(t\in\mathbb R^2\):
\[
 E e^{it\cdot W/\sqrt n}
       =1-\frac{t^{\mathsf T}\operatorname{Cov}(W)t}{2n}
          +O(n^{-3/2}).
\]
The product for the independent summands tends to
\(\exp[-t^{\mathsf T}\operatorname{Cov}(W)t/2]\).
One can pass from this calculation to bounded-continuous-test
convergence as follows. Add an independent \(N(0,\delta^2 I_2)\).
Its integrable Gaussian Fourier factor permits Fubini and the
one-dimensional Gaussian integral in each coordinate; dominated
convergence then gives uniform convergence of the smoothed
densities. Uniform second-moment bounds give tightness, so density
convergence on a large compact set gives convergence of bounded
continuous tests. Finally let \(\delta\downarrow0\): uniform
continuity on that compact set and tightness bound the change
caused by the added Gaussian, uniformly in \(n\).
This proves the required bivariate limit, including a zero
second variance. Its factorized Gaussian characteristic function
gives independent coordinates.

At the empirical root \(m_n(H_n/\sqrt n)=0\).
Tightness lets the uniform local expansion be evaluated there,
so \(\mathcal D(H_n)=Z_n+o_P(1)\).
The inverse Lipschitz bound implies
\[
                   H_n=\mathcal D^{-1}(Z_n)+o_P(1).
\]
Continuity of the inverse and bounded-continuous-test convergence,
followed by uniform continuity on compact sets for the vanishing
perturbation, give the asserted law. Every atom at the join has
remained in both the population and empirical expansions.
\end{proof}
\end{annexitem}

\begin{annexitem}{annex:p15.huber.joint.law.variance}{thm:p15-huber-joint-law-and-variance}{Gaussian special case, nonsmooth law and the proposed variance estimate}
\begin{proof}
When \(p=0\), the inverse is the diagonal linear map with entries
\(1/A,1/C\), proving the Gaussian special case.
When \(p>0\), \(\psi(U)^2\) equals \(k^2\) with positive probability,
while its mean \(\beta\) is strictly smaller than \(k^2\).
It is not constant, so \(\eta^2>0\).
The driving Gaussian then has a positive density everywhere.
On an open cone, change of variables \(z=M_eh\) gives the
displayed density, with positive determinant.
The cones where both positive parts vanish and where both
are active have determinants
\[
                AC,\qquad (A+2p)(C+2pk^2),
\]
respectively. These are different.
If the joint law were Gaussian it would be nondegenerate, because
its density is positive on every open cone. Its density would
be continuous at zero. Equality almost everywhere with the
displayed analytic density on an open cone implies equality
everywhere within that cone, by continuity.
Approaching zero within the two cones would then require the
same limiting density value, contrary to the different determinants.
This proves joint non-Gaussianity.

For \(a\geq0\), the positive-part difference in \(\mathcal D_1\)
lies between zero and \(2a\). Hence
\[
                       Aa\leq z_1\leq(A+2p)a.
\]
Reflection handles negative \(a\), and \(a=0\) gives \(z_1=0\).
It follows that
\[
 |Z_1|/(A+2p)\leq |H_1|\leq |Z_1|/A .
\]
The upper and lower cones each have positive probability when
\(p>0\); on suitable open subsets each inequality is strict.
Squaring and integrating proves the strict moment bounds.
The reflection \((a,b)\mapsto(-a,b)\) reflects only the first
coordinate of \(\mathcal D\). The driving Gaussian has this
symmetry, so \(H_1\) has mean zero; integrability follows from
the Lipschitz inverse. Weak convergence applied to the bounded
continuous function \(\min(S_0^2h_1^2,R^2)\), followed by
monotone convergence in \(R\), proves the ordered truncated-risk
identity. It proves no untruncated uniform integrability assertion.

For the variance-estimate consequence, a boundary observation
\(U=k\) is strictly interior at the estimated parameters exactly
when \(H_{n,1}+kH_{n,2}>0\), once the estimates are close enough
that the opposite boundary is irrelevant.
At \(U=-k\) the corresponding condition is
\(-H_{n,1}+kH_{n,2}>0\).
The empirical mass near but not at the two boundaries tends
to zero by the band argument in the limit proof.
The strictly interior empirical fraction consequently equals
\[
 A+p{\bf1}_{\{H_{n,1}+kH_{n,2}>0\}}
   +p{\bf1}_{\{-H_{n,1}+kH_{n,2}>0\}}+o_P(1).
\]
The limiting law gives no mass to the two boundary lines,
because their images are lines under the piecewise nonsingular
linear map and \(Z\) has a density.
Approximating the indicators by continuous functions outside
arbitrarily thin strips therefore permits passage to their
limiting law.
At a root the squared-score average is exactly \(\beta\);
the scale converges to \(S_0\), and \(n/(n-1)\to1\).
Substitution gives the asserted random limit.
The lower, side and upper cones have positive probability,
giving all three different denominator values.
A sequence converging in probability to a constant would have
that constant as its distributional limit, so no such convergence
is possible here.
\end{proof}
\end{annexitem}

\begin{annexitem}{annex:p15.huber.joint.worked.atomic.law}{thm:p15-huber-joint-worked-atomic-law}{An exactly calculated four-point law}
\begin{proof}
At \(S_0=1\), the squared score takes values \(1/4\) and one,
each with probability \(1/2\). Thus its mean is \(5/8\) and
its fourth moment is \(17/32\), giving
\(\eta^2=(17/32-25/64)/4=9/256\).
The interior mass and second moment are \(1/2,1/8\).
The inverse matrices for the lower, upper and right cones are,
respectively,
\[
 \begin{pmatrix}2&0\\0&8\end{pmatrix},\qquad
 \begin{pmatrix}1&0\\0&8/5\end{pmatrix},\qquad
 \frac17\begin{pmatrix}12&-8\\-8&24\end{pmatrix}.
\]
The left cone is the reflected right one.
Substituting their candidate coordinates into the cone inequalities
gives exactly \(w\leq-r/4\), \(w\geq5r/8\), and the intervening
region. This proves the entire piecewise formula, including its
scale coordinate and matching boundary values.
The three variance-estimate limits are
\((5/8)/(1/2)^2\), \((5/8)/(3/4)^2\), and \(5/8\).

Here is an exact evaluation of the location variance.
Let
\[
 s_1=\sqrt{5/8},\quad s_2=3/16,\quad
 \theta_-=-\arctan(\sqrt{10}/3),\quad
 \theta_+=\arctan(5\sqrt{10}/6),
\]
\[
 F_c(t)=t/2+\sin(2t)/4,\quad
 \alpha=12s_1/7,\quad d=8s_2/7,
\]
and
\[
 J(t)=\tfrac12(\alpha^2+d^2)t
       +\tfrac14(\alpha^2-d^2)\sin(2t)
       +\tfrac12\alpha d\cos(2t).
\]
Then \(F_c'=\cos^2t\) and
\(J'=(\alpha\cos t-d\sin t)^2\).
Write the two independent standard normal coordinates in polar
form. Their angle is uniform, their radius is independent with
second moment two, as follows by the polar change of variables
in their joint density. Restricting to positive first coordinate
and using reflection, the piecewise formula gives
\[
 EH_1^2=\frac2\pi\left[
 4s_1^2\{F_c(\theta_-)-F_c(-\pi/2)\}
 +J(\theta_+)-J(\theta_-)
 +s_1^2\{F_c(\pi/2)-F_c(\theta_+)\}\right].
\]
Evaluation gives \(1.832700020\ldots\).
This is an integral calculation of the limit distribution,
not a sampling estimate.
\end{proof}
\end{annexitem}

\begin{annexitem}{annex:p15.huber.joint.normal.contamination.counterexample}{thm:p15-huber-joint-normal-contamination-counterexample}{A non-Gaussian location limit inside the original model}
\begin{proof}
We first establish a sufficient condition making the location
marginal non-Gaussian, and then verify it in the specified
normal-contamination law.
In this proof \(k=1\), \(p>0\), and
\[
 \Delta=AC+p(A+C),\quad
 \alpha=(C+p)/\Delta,\quad d=p/\Delta,\quad
 \ell=1/(A+2p),\quad u=1/A.
\]
The cone inverse gives, with \(R=|Z_1|\) and \(W=Z_2\),
\[
 |H_1|=\min\{uR,\max(\ell R,\alpha R-dW)\},
 \qquad \operatorname{sign}H_1=\operatorname{sign}Z_1.
\]
The lower driver boundary is \(W=-(C/A)R\), and the upper one
is \(W=(C+2p)R/(A+2p)\).
Since \(A-C=E[(1-U^2){\bf1}_{\{|U|<1\}}]>0\),
\[
 \alpha-\ell=\frac{p(C+2p)}{\Delta(A+2p)}
          >\frac{pC}{\Delta A}=u-\alpha>0 .
\]
Pair the equally likely values \(W=w\) and \(W=-w\) conditional
on \(R>0\). Until \(dw=(u-\alpha)R\), neither value is clipped.
Until \(dw=(\alpha-\ell)R\), only the larger value is clipped,
strictly decreasing their sum of squares.
Beyond that second threshold both are clipped. At the threshold
the unclipped pair is \(\ell R,(2\alpha-\ell)R\), with
\(2\alpha-\ell>u\); its sum of squares already exceeds the
clipped sum \((\ell^2+u^2)R^2\), and the unclipped sum increases
with \(w^2\) afterward.
The normal \(W\) gives positive probability to these strict regions.
Consequently
\[
              EH_1^2<\alpha^2\beta+d^2\eta^2=:v.
\]

Suppose now that
\[
                     pA\eta^2<C(C+p)\beta.                 \tag{*}
\]
For the Gaussian linear form \(L=\alpha Z_1-dZ_2\), the
direction of its most likely large positive values has coordinate
ratio
\[
       Z_2/Z_1=-d\eta^2/(\alpha\beta)
               =-p\eta^2/((C+p)\beta).
\]
Condition (*) places this strictly inside the right middle cone.
We justify the tail conclusion explicitly. Whiten the two
Gaussian coordinates, and rotate their standard-normal vector
so that its first coordinate is \(L/\sqrt v\).
The rotated coordinates remain independent standard normals,
by their product density under an orthogonal transformation.
In that representation the mean direction along the first
coordinate lies strictly inside the cone. There is therefore
a constant \(c>0\) such that the middle-cone condition holds
whenever the first coordinate is positive and the absolute
second coordinate is at most \(c\) times the first.
Conditioning on \(L>x\) makes the probability of the contrary
event at most \(P(|N(0,1)|>cx/\sqrt v)\), tending to zero.
It follows that
\[
 P(L>x,\ \hbox{right middle cone})/P(L>x)\longrightarrow1.
\]

For \(Z_1>0\), downward clipping only reduces \(L\).
The only way that \(H_1>x\) is not covered by \(L>x\) is the
upper cone, where \(H_1=\ell Z_1\). That event implies
\[
 Z_1>(A+2p)x,\qquad Z_2>(C+2p)x .
\]
Its probability, by independence and normal tail bounds, decays
at exponential rate
\(\frac12[(A+2p)^2/\beta+(C+2p)^2/\eta^2]x^2\).
This rate is strictly larger than \(x^2/(2v)\).
Indeed the vector \((A+2p,C+2p)\) has \(L\)-value one;
Cauchy--Schwarz in the covariance metric gives squared norm
at least \(1/v\), with equality only in the covariance-weighted
direction of \((\alpha,-d)\), whose second coordinate is negative.
The displayed vector has positive second coordinate, so equality
is impossible.
For precision, the elementary normal bounds
\(\phi(t)/(t+t^{-1})\leq1-\Phi(t)\leq\phi(t)/t\), \(t>0\),
follow by integrating \(\phi(x)/x\) by parts and bounding the
remaining \(x^{-2}\) term. They justify both the rate and the
negligibility relative to \(P(L>x)\).
Combining the lower and upper estimates gives
\[
                    P(H_1>x)/P(N(0,v)>x)\longrightarrow1.
\]
The location law is symmetric and has a finite second moment.
If it were Gaussian, this tail ratio would force its variance
to be \(v\), contradicting the strict second-moment inequality.
This proves the sufficient condition.

We now verify every premise for the original-model example.
The function \(B(q)/q^2\) is continuous and strictly decreasing
from one to zero. The proposed right side is positive and
smaller than \(B(1)\), so its unique root exceeds one.
To bound it above by two, first note
\[
 B(1)=2\int_0^1x^2\phi(x)\,dx+2(1-\Phi(1))
            \geq(2/3)\phi(1)+\phi(1)>1/3 .
\]
The normal tail inequality just proved gives the middle bound,
and \(\phi(1)>1/5\) follows from \(e<3,\pi<4\).
Meanwhile \(B(2)/4<1/4\), since \(B(2)<EX^2=1\).
Thus \((B(1)-2/1000)/(1-2/1000)>1/4\), proving \(q<2\).
At \(S_0=q\), the normal part contributes
\(a_0 B(q)/q^2\) to the squared-score mean, and the two atoms
contribute \(2p\). Their sum is exactly \(\beta=B(1)\).
The law has no atom at zero, so it is identified.

The interior constants are those in the statement. They satisfy
\[
 A<1,\qquad
 C\geq\frac{a_0}{4}\,2\int_0^1x^2\phi(x)\,dx
       >a_0/30>1/40.
\]
Since \(0\leq\psi(U)^2\leq1\),
\(\eta^2\leq\beta/4\). Hence
\[
 pA\eta^2<\beta/4000
           <\beta/1600<C(C+p)\beta ,
\]
verifying (*).
Both coordinates of the driving Gaussian are nondegenerate
because \(p>0\). The complete non-Gaussian marginal argument
therefore applies inside this normal-contamination model.
The decimal constants in the body evaluate the unique scalar
root and its displayed integrals.
\end{proof}
\end{annexitem}

\begin{annexitem}{annex:p15.huber.joint.atomic.bound}{thm:p15-huber-joint-atomic-bound}{The sharp source upper bound survives with its corrected interpretation}
\begin{proof}
The directional-map inequality already gives
\(S_0^2EH_1^2\leq S_0^2\beta/A^2\).
Put \(q=kS_0\). For an identified contaminated normal law,
\[
 S_0^2\beta=q^2r
   =(1-\varepsilon)B(q)+\varepsilon E_{H_{\rm c}}\psi_q(X)^2
   \leq(1-\varepsilon)B(q)+\varepsilon q^2,
\]
where \(H_{\rm c}\) denotes the contaminating law here.
This algebra uses no derivative assumption or exclusion of atoms.
As in the earlier calculation, it implies \(q\leq q_{\max}\).
The strict interior mass obeys
\(A\geq(1-\varepsilon)A(q)\), and \(A(q)/q\) decreases in \(q\).
Therefore
\[
 S_0^2EH_1^2
 \leq\frac{q^2r}{(1-\varepsilon)^2A(q)^2}
 \leq\frac{q_{\max}^2r}{(1-\varepsilon)^2A(q_{\max})^2}
 =W(q_{\max},\varepsilon).
\]
Every symmetric contaminant strictly outside the final cutoff
has \(q=q_{\max}\), no cutoff atom, and equality throughout.
It is identified since it gives no mass to zero.
This proves sharpness over identified laws.
If \(\varepsilon<1-r\), every contaminant gives
\(F\{0\}\leq\varepsilon<1-r\); thus none are removed by the
identification condition. This establishes the asserted
whole-class quantifier and the corrected interpretation of the
historical upper bound.
\end{proof}
\end{annexitem}

\begin{annexitem}{annex:p15.huber.zero.atom.mixture.bridge}{thm:p15-huber-zero-atom-mixture-bridge}{The normal-mixture comparison needs no contamination-fraction restriction}
\begin{proof}
If \(F\{0\}=0\), then \(0<\beta<k^2\) implies the strict
identification inequality immediately. The identification proof
already gives eventual unique estimation and almost-sure
consistency without any condition on the largest atom.
For a non-atomic law the two clipping atoms have mass zero,
so the joint limit theorem gives its linear Gaussian case.

Each displayed normal mixture has a positive continuous
symmetric density, including its two endpoint mixtures.
Its mass at zero and its masses at every possible cutoff are zero.
For \(k=1.5\), \(0<B(k)<k^2\).
The population scale equation and variance therefore have
exactly the formulas in the predecessor chapter:
the scale-three component contributes \(9B(q/3)\) to the
unnormalized squared-score moment and \(A(q/3)\) to its interior
probability. The uniquely defined scale and the new limit theorem
justify them for all the mixture fractions stated here.
No constraint on that fraction has been inferred from a worst-case
bound on atoms in a different family.
\end{proof}
\end{annexitem}

\begin{annexitem}{annex:p15.huber.scale.sign.minimax}{thm:p15-huber-scale-sign-minimax}{The finite sign-conditioning lower bound and full minimax equality}
\begin{proof}
The observation-domain convention changes no possible behavior
on this model. Every competing function on vectors with all
coordinates nonzero extends measurably and positively to vectors
with some zero coordinates but not all zero: use
\(\max_i|x_i|\) on those additional vectors. Scaling preserves
that set, so the extension remains scale equivariant.
Conversely any such larger-domain estimator restricts to the stated
class. All source laws give the additional vectors probability zero.

Write \(a_0=1-\varepsilon\) and \(c=(q^2-1)/2\).
An explicit least-favorable density on the original observation
space is
\[
 f_*(x)=
 \begin{cases}
 a_0\phi(x),&|x|\leq q,\\
 a_0\phi(q)(|x|/q)^{-q^2},&|x|>q.
 \end{cases}
\]
Its integral is
\(a_0[A(q)+2q\phi(q)/(q^2-1)]=1\), by the exact scale
normalization. Outside the central interval its ratio to
\(a_0\phi(x)\) has logarithm
\[
       (x^2-q^2)/2-q^2\log(|x|/q)\geq0;
\]
on the positive half-line its derivative is \(x-q^2/x\geq0\)
and the value at \(q\) is zero. Reflection gives the negative
half-line. Thus \(f_*-a_0\phi\) is a nonnegative density of
mass \(\varepsilon\), supported strictly outside the cutoff
up to null endpoints. This proves \(F_*\in\mathcal F_q\).

Its even density makes the sign a fair independent sign,
independent of \(Y=\log X^2\). The density of \(Y\) is
\(g_*(y)=e^{y/2}f_*(e^{y/2})\).
Below \(y=\log q^2\) it is \(a_0g_Y(y)\); above it,
\[
                 g_*(y)=a_0q\phi(q)e^{-c(y-\log q^2)}.
\]
This is exactly the positive \(C^1\) least-favorable transformed
density with finite information \(I_Y\) from the scale construction.
For an i.i.d.\ sample, each sign pattern
\(s=(s_1,\ldots,s_n)\in\{-1,1\}^n\) has probability \(2^{-n}\);
conditional on it the \(Y_i\) remain independent with density \(g_*\).

Fix any competing scale estimator, any \(n\), and any sign pattern.
Define the real measurable function
\[
 T_{n,s}(y_1,\ldots,y_n)
       =\log S_n(s_1e^{y_1/2},\ldots,s_ne^{y_n/2})^2 .
\]
It is finite because the original estimator is finite and positive.
For every \(t\in\mathbb R\), scale equivariance with multiplier
\(e^{t/2}\) gives
\(T_{n,s}(y+t{\bf1})=T_{n,s}(y)+t\).
This is the original estimator evaluated conditional on known signs,
not a newly randomized estimator or an averaged decision.
The finite integrated-score bound
\cref{thm:p15-location-integrated-score-bound} therefore gives,
for every \(R>0\),
\[
 E_{g_*}\min(nT_{n,s}^2,R^2)
                     \geq (I_Y+4\pi^2/R^2)^{-1}.
\]
Sum this inequality with the actual weights \(2^{-n}\).
Its right side is the same for every pattern. The resulting
left side is precisely the expectation under the original \(F_*\)
in the statement. The bound is finite-sample and uniform over
all the conditional functions, so patterns may vary with \(n\);
no limiting exchange over sign patterns is needed.

Taking \(\liminf_n\), then letting \(R\to\infty\), then using
the supremum over laws proves the lower bound \(1/I_Y\) for
every estimator sequence in the original comparison class.
For the upper bound, the magnitude estimator solves
\[
                    n^{-1}\sum_i\min(X_i^2/S^2,q^2)=1.
\]
For every vector with all coordinates nonzero this continuous monotone
equation has exactly one positive root: its limiting values
are \(q^2>1\) and zero, and an unsaturated nonzero observation
makes the decrease strict at the target. Its root is measurable,
since its comparison with a fixed positive number is determined
by the corresponding continuous score inequality. Substitution
shows scale equivariance directly.
For every \(F\in\mathcal F_q\), the transformed contamination
has constant score \(c\) on its strict tail, independently of
the original signs. The scale theorem gives
\(\sqrt n\log S_n^2\Rightarrow N(0,1/I_Y)\).
Bounded continuous truncated losses converge, and increasing
the truncation afterward gives risk \(1/I_Y\) for every such \(F\).
Finally, pathwise equivariance gives
\[
 \log S_n(\sigma X)^2-\log\sigma^2=\log S_n(X)^2.
\]
The supremum over \(\sigma\) changes no risk, even before limits.
The lower and upper bounds therefore have precisely the stated
quantifiers and limiting order, proving equality for all positive
scale-equivariant estimators allowed to use signs.
\end{proof}
\end{annexitem}
